\documentclass{minmal}
\begin{document}
 \begin{center}
    \textbf{Prøve om syre og base:}
\end{center}
\textbf{Oppg. 1:}\nl[1]
a) Hva menes med en syre og hva menes med en base?\nl[1]
\textbf{svar:} En syre er et stoff som avgir protoner ($H^+$). En base er et stoff tar opp protoner ($H^+$).\nl[1]
b) Forklar pH skalaen (stikkord: basisk,sur, nøytral, hvilke verdier har vi).\nl[1]
\textbf{Svar:} pH skalaen går fra 0 til 14.  0 til 7 er surt, 7 til 14 er basisk. ved pH lik 7 er det en nøytral løsning. Da er $[H_3O^+]=[OH^-]=10^{-7}M$\nl[1]
c) Hvordan defineres pH?\nl[1]
\textbf{Svar:} pH er en matematisk definisjon:\nl[1]
$$pH=-log([H_3O^+])$$\nl[1]
d) Skriv reaksjonen med vann for følgende syrer:$HCl,\ HNO_3,\ H_2SO_4$.\nl[1]
$$HCl+H_2O=Cl^-+H_3O^+$$
$$HNO_3+H_2O=NO_3^-+H_3O^+$$
$$H_2SO_4+2H_2O=SO_4^{2-}+2H_3O^+$$  eller:\\
$$H_2SO_4+H_2O=HSO_4^{-}+H_3O^+$$
e) Skriv reaksjonen som skjer når følgende salt løses i vann: $NaOH,\ Ca(OH)_2$\nl[1]
$$NaOH=Na^++OH-$$
$$Ca(OH)_2=Ca^{2+}+2OH^-$$
\newpage
f) Fullfør følgende reaksjon:\nl[1]
$$CaO + H_2O =$$\nl[1]
$$CaO+H_2O=Ca^{2+}+2OH^-$$\\
\textbf{Oppg. 2:}\nl[1]
a) Hva er $[H_3O+]$ når pH=3,5?\nl[1]
\textbf{Svar:}  Her gjelder følgende formel: \\
$$[H_3O^+]=10^{-pH}=10^{-3,45}=0,00032$$
b) Beregn pH i følgende syrer: $0,4M\ HNO_3\ 0,1M HCl\ 0,1M H_3PO_4$
(anta at $H_3PO_4$ protlyserer fullstendig).\nl[1]
\textbf{Svar:} Vi starter med raksjonoen:
$$HNO_3+H_2O=NO_3^-+H_3O^+$$
Her vil $[H_3O^+]=[HNO_3]=0,4M$ Vi får $pH=-log(0,4)=0,4$
\nl[1]
$$HCl+H_2O=Cl^-+H_3O^+$$
Her vil $[H_3O^+]=[HCl]=0,1M$ Vi får $pH=-log(0,1)=1$
\nl[1]
$$H_3PO_4+3H_2O=PO_4^{3-}+3H_3O^+$$
Her er det viktig å vite at fosforsyre ($H_3PO_4$) er en treprotisk syre. Vi må ha 3 vannmolekyler og får da dannet 3 $H_3O^+$. \\ 
Det betyr at:$$[H_3O^+]=3\cdot [H_3PO_4]=3\cdot 0,1M=0,3M$$\\
For pH får vi da: $pH=-log(0,3)=0,52$\newpage 
\noindent c) Beregn pH i følgende baser: 0,1M NaOH, 0,1M $Ca(OH)_2$,\ 0,2M CaO\nl[1]
\textbf{0,1M NaOH:}  NaOH er et salt som løser seg i vann etter følgende reaksjon:
$$NaOH=Na^++OH^-$$
Det betyr at $[OH^-]=[NaOH]=0,1M$
Her er det $[OH^-]$ vi kjenner, og da må vi beregne $pOH$:
$$pOH=-log([OH^-])=-log(0,1)=1$$
Videre  så er $pH=14-pOH=14-1=13$\\
\textbf{0,1M $\mathbf{Ca(OH)_2}$:}  Her vil saltet $Ca(OH)_2$ løses i vann:
$$Ca(OH)_2=Ca^{2+}+2OH^-$$
Her er det kosnetrasjonen av $[OH^-]$ som er interessant.\\
Her vil $[OH^-]=2\cdot[Ca(OH)_2]$ for vi ser at vi har dobbelt så mange $OH^-$ som $Ca(OH)_2$. Vi får:
$$[OH^-]=2\cdot [Ca(OH)_2]=2\cdot 0,1M=0,2M $$
$$pOH=-log([OH^-])=-log(0,2)=0,7$$
$$pH=14-pOH=14-0,7=13,3$$
\textbf{0,2M CaO:}\\
Her gjelder følgende reaksjon:
$$CaO+H_2O=Ca^{2+}+2OH^-$$
Det betyr at $[OH^-]=2\cdot [CaO]=2\cdot 0,2M=0,4M$
Vi får:
$$pOH=-log([OH^-])=-log(0,4)=0,4$$
$$pH=14-pOH=14-0,4=13,6$$
d) Hva blir [NaOH] i følgende titreringer?\nl[1]
Her må vi bruke fortynningsloven som gjelder for syre-base titrering.  Den sier at det ved ekvivalenspunktet er like mange mol av syre og base:\nl[1]
$$C_b\cdot V_b=C_s\cdot V_s$$
b står for base, s for syre. Det er $NaOH$ som er basen i dette tilfellet.  Vi løser ligningen m.h.p.  $C_b$:
$$C_b=\frac{C_s\cdot V_s}{V_b}$$
\textbf{Kommentar:} Volumet skal gis i liter. Men hvis begge volumene er gitt i mL så kan man bruke det.\newpage  
Vi får:\\\\
1) 25mL 0,1M HCl med 20mL NaOH\nl[1]
$$C_{NaOH}=\frac{C_{HCl}\cdot V_{HCl}}{V_{NaOH}}=\frac{0,1M\cdot 25mL}{20mL}=0,125M$$
2) 15mL 0,2M HCl med 30mL NaOH\nl[1]
$$C_{NaOH}=\frac{C_{HCl}\cdot V_{HCl}}{V_{NaOH}}=\frac{0,2M\cdot 15mL}{30mL}=0,1M$$
3) 30mL 0,3M HCl med 30mL NaOH\nl[1]
$$C_{NaOH}=\frac{C_{HCl}\cdot V_{HCl}}{V_{NaOH}}=\frac{0,3M\cdot 30mL}{30mL}=0,3M$$
4) 25mL 0,25M HCl med 20mL NaOH\nl[1]
$$C_{NaOH}=\frac{C_{HCl}\cdot V_{HCl}}{V_{NaOH}}=\frac{0,25M\cdot 25mL}{20mL}=0,313M$$
\end{document}